\documentclass{article}
\usepackage[utf8]{inputenc}
\usepackage{a4wide}

\begin{document}

\section*{ไปให้ไกล ๆ}

ชลธเนติ์มีกองกระดาษเหลือใช้จำนวนหนึ่ง

เธอต้องการใช้เศษกระดาษเหล่านี้ที่มีทั้งหมดมาจำลองการสร้างถนน

เธอจึงนำเศษกระดาษทั้งหมดมาม้วนเป็นแท่งเส้นตรงที่มีขนาดความยาวไม่เท่ากันขึ้นกับเศษกระดาษที่มีในแต่ละใบ

เพื่อให้ง่ายในการสร้าง ชลธเนติ์จำลองจุดเริ่มต้นของถนนที่พิกัด (0,0)

ม้วนกระดาษแต่ละอันที่เอามาต่อสามารถตั้งอยู่ในแนวนอน (ขนานกับแกน X) หรือแนวตั้ง

(ขนานกับแกน Y) โดยที่เมื่อเธอต่อม้วนกระดาษม้วนใหม่

จะต้องไม่ต่อกับม้วนกระดาษอันล่าสุดในทิศทางเดียวกัน (ถ้าอันหนึ่งแนวนอน อีกอันต้องตั้ง

หรือถ้าอันหนึ่งตั้ง อีกอันต้องแนวนอน)



งานของคุณ: ช่วยชลธเนติ์หาค่าระยะห่างที่มากที่สุดจากจุดเริ่มต้น

ถึงจุดปลายม้วนกระดาษม้วนสุดท้าย ซึ่งจะต้องหาค่าระยะห่างในรูปของ ระยะทางยกกำลังสอง

(ไม่ต้องคำนวณเป็นระยะทางจริง)



\textbf{Input:}



บรรทัดแรกเป็นจำนวนเต็มแทนจำนวนม้วนกระดาษจำนวน n อัน (1 ≤ n ≤ 100000)



บรรทัดที่ 2 ความยาวของม้วนกระดาษแต่ละอันเป็นจำนวนเต็ม a1, a2, ..., an โดยที่ 1 ≤

ai ≤ 10000



\textbf{ตัวอย่าง}



\textbf{Input:}



3



5 2 3



\textbf{Output:}



68



\textbf{Input:}



4



1 2 3 3



\textbf{Output:}



45



\subsection*{Input}
บรรทัดแรก : เป็นจำนวนเต็มแทนจำนวนม้วนกระดาษจำนวน n อัน (1 ≤ n ≤ 100000)



บรรทัดที่ 2 : ความยาวของม้วนกระดาษแต่ละอันเป็นจำนวนเต็ม a1, a2, ..., an โดยที่ 1 ≤

ai ≤ 10000



\subsection*{Output}
{ค่าระยะห่างที่มากที่สุดจากจุดเริ่มต้น ถึงจุดปลายม้วนกระดาษม้วนสุดท้ายซึ่ง{ในรูปของ

ระยะทางยกกำลังสอง}}\\



\end{document}
